\documentclass[11pt]{article}
\usepackage[T1]{fontenc}
\usepackage{lmodern}
\usepackage[margin=1in]{geometry}
\usepackage{amsmath,amssymb,amsthm}
\usepackage[hidelinks]{hyperref}
\newtheorem{theorem}{Theorem}
\newtheorem{lemma}{Lemma}
\newcommand{\R}{\mathbb R}
\hypersetup{pdftitle={A counterexample to exact boundary control for a parabolic obstacle problem},pdfauthor={Kaiwen Jin}}
\title{A counterexample to exact boundary control\newline for a parabolic obstacle problem}
\author{Kaiwen Jin\\\small\texttt{yui.ui.siki@gmail.com}}
\date{October 12, 2026}

\begin{document}
\maketitle

\begin{abstract}
We give a smooth one-dimensional obstacle and a smooth admissible initial
state for which no endpoint controls can steer the parabolic obstacle
problem exactly to its stationary energy minimizer in finite time.
The obstacle has two interior maxima of equal height. The stationary
solution is constant between them, while the obstacle prevents the evolving
state from falling below that height at either end of this interval.
A positive sine test function then gives a strictly positive lower bound on
the terminal $L^2$ error for every finite time. This answers the universal
boundary-control question in AIM Problem 398 negatively.
\end{abstract}

\section{The control question}

Let $I=(0,1)$ and let $\psi\in C^\infty([0,1])$ have negative endpoint
values and a positive maximum. Write $\bar y$ for the unique minimizer of
\[
 E(u)=\int_I |u_x|^2\,dx
 \quad\text{over}\quad
 K=\{u\in H^1_0(I):u\ge\psi\text{ a.e.}\}.
\]
The parabolic obstacle problem evolves above the same lower obstacle:
\begin{equation}\label{eq:obstacle}
 \min\{y_t-y_{xx},y-\psi\}=0,\qquad
 y(t,i)=g_i(t)\ (i=0,1),\qquad y(0)=y_0.
\end{equation}
The equation is interpreted by the associated parabolic variational
inequality. In particular, the state stays above $\psi$, and its heat
operator is nonnegative in the weak sense.

Boundary control to the stationary obstacle solution is discussed by
Geshkovski \cite[Section 1.5.1, pp.~37--39]{geshkovski}.
Pighin and Zuazua \cite[Section 7.3]{pighin} explain why controllability of
penalized semilinear equations does not by itself give a control result for
the obstacle limit. The concrete assertion in AIM Problem 398
\cite{aim398} allows every smooth obstacle of the above kind. It asks whether
every $y_0\in K$ admits some finite $T>0$ and controls
$g_0,g_1\in H^1(0,T)$ satisfying
\begin{equation}\label{eq:controls}
 g_i(0)=g_i(T)=0,\qquad g_i(t)\ge\psi(i),
\end{equation}
for which $y(T)=\bar y$ in $L^2(I)$.

Two separated equal maxima give a counterexample to this assertion.
Between them the stationary state is a horizontal segment, and the two
internal obstacle values prevent the controls from removing a positive heat
moment in finite time. Only the state constraint at these fixed points is
needed; they need not remain contact points of the evolving solution.

All functions below are real-valued. Spatial point values of $H^1$ functions
use their continuous representatives on the closed interval
\cite[Propositions 1.8.2, 1.8.5 and 1.8.7]{gerard}.

\section{The counterexample}

\begin{theorem}\label{thm:counterexample}
Set
\begin{equation}\label{eq:data}
 \psi(x)=1-81\left(x-\tfrac13\right)^2\left(x-\tfrac23\right)^2,
 \qquad y_0(x)=9x(1-x),\qquad x\in[0,1].
\end{equation}
Then $\psi(0)=\psi(1)=-3$, $\max\psi=1$, and
$y_0\in H^1_0(I)$ with $y_0\ge\psi$.
Let $\bar y$ be the minimizer defined above. For every finite $T>0$ and every
pair of controls $g_0,g_1\in H^1(0,T)$ satisfying \eqref{eq:controls},
the variational solution of \eqref{eq:obstacle} satisfies
\begin{equation}\label{eq:error}
 \|y(T)-\bar y\|_{L^2(I)}
 \ge \frac{2\sqrt6}{3\pi}\,e^{-9\pi^2T}>0.
\end{equation}
In particular, no such controls reach $\bar y$ at any finite time.
\end{theorem}

We first identify the part of the stationary state used in the proof.
\begin{lemma}\label{lem:plateau}
For the obstacle in \eqref{eq:data}, $\bar y(x)=1$ on
$[\tfrac13,\tfrac23]$.
\end{lemma}
\begin{proof}
The obstacle is at most $1$ and equals $1$ at $a=\tfrac13$ and $b=\tfrac23$.
The function $\min\{\bar y,1\}$ is an admissible competitor with no greater
energy: Sobolev truncation gives
$|(\min\{\bar y,1\})_x|\le|\bar y_x|$ a.e.
\cite[Proposition 3.22]{hunter}. For example, apply the weak derivative
formula for absolute values to
$\min\{\bar y,1\}=(\bar y+1-|\bar y-1|)/2$.
Uniqueness of the minimizer therefore gives $\bar y\le1$.
Continuity and the obstacle constraint imply $\bar y(a)=\bar y(b)=1$.

Replace $\bar y$ by $1$ on $[a,b]$, leaving it unchanged elsewhere.
Matching the values at $a,b$ makes this an $H^1_0(I)$ competitor; its weak
derivative equals $\bar y_x$ outside $(a,b)$ and vanishes inside.
It remains above $\psi$, since $\psi\le1$. Minimality gives
\[
 E(\bar y)\le E(\bar y)-\int_a^b|\bar y_x|^2\,dx.
\]
Hence $\bar y_x=0$ a.e. on $(a,b)$, and its endpoint value gives the result.
\end{proof}

\begin{proof}[Proof of Theorem~\ref{thm:counterexample}]
Both functions in \eqref{eq:data} are polynomials, and $y_0$ has zero
endpoint values. With $r=x(1-x)$, direct expansion yields
\[
 y_0-\psi=9r-1+81(\tfrac29-r)^2
 =81(r-\tfrac16)^2+\tfrac34>0.
\]
The square in the definition of $\psi$ also shows that its maximum is $1$
and its endpoint values are $-3$. Thus all the data conditions hold.

Fix any finite $T$ and any controls allowed in the theorem. We use the
energy interpretation of their variational solution:
\begin{equation}\label{eq:energy}
 y\in C([0,T];L^2(I))\cap L^2(0,T;H^1(I)),\qquad y\ge\psi\text{ a.e.}
\end{equation}
The boundary traces are imposed through the affine lift
$\ell(t,x)=(1-x)g_0(t)+xg_1(t)$, with
$y-\ell\in L^2(0,T;H^1_0(I))$. The given $H^1$ controls make
$\ell_t\in L^2(0,T;L^2(I))$.
For every nonnegative $\phi\in H^1_0(I)$, the upward variation
$v=y+s\phi$, $s\ge0$, has the same boundary traces and is admissible.
Its variational inequality, integrated against a nonnegative
$\eta\in C^\infty_c(0,T)$, gives
\begin{equation}\label{eq:upward}
 -\int_0^T\eta'(t)\int_I y\phi\,dx\,dt
 +\int_0^T\eta(t)\int_I y_x\phi_x\,dx\,dt\ge0.
\end{equation}
This formulation uses a distributional time derivative and requires no
classical regularity of the contact boundary.

Let $a=\tfrac13$, $b=\tfrac23$, and $k=3\pi$. To measure the deviation on
the stationary plateau, take the nonnegative function
\[
 \varphi(x)=
 \begin{cases}
  \sin(k(x-a)),&a<x<b,\\
  0,&x\notin(a,b).
 \end{cases}
\]
It is continuous and piecewise smooth, with weak derivative
$k\cos(k(x-a))$ on $(a,b)$ and zero elsewhere, so
$\varphi\in H^1_0(I)$. At almost every time, spatial continuity and the
obstacle constraint give $y(t,a),y(t,b)\ge1$.
These traces belong to $L^2(0,T)$ by the bounded point-evaluation maps on
$H^1(I)$.

Define the continuous moment $m(t)=\int_a^b y(t,x)\varphi(x)\,dx$.
Integrating in space on $(a,b)$, where the sine is smooth, gives
\begin{align*}
 \int_I y_x\varphi_x\,dx
 &= [y\varphi_x]_a^b-\int_a^b y\varphi_{xx}\,dx\\
 &= -k\{y(t,a)+y(t,b)\}+k^2m(t).
\end{align*}
The boundary terms use $\varphi_x(a+)=k$ and $\varphi_x(b-)=-k$.
Thus \eqref{eq:upward} implies, in time distributions,
\[
 m'+k^2m\ge k\{y(t,a)+y(t,b)\}\ge2k.
\]
Since $\int_a^b\varphi\,dx=2/k$, the deviation moment
$F(t)=m(t)-2/k$ satisfies $F'+k^2F\ge0$.

It follows that $H(t)=e^{k^2t}F(t)$ is nondecreasing. To justify this for
the distributional inequality, mollify the continuous function $H$ on any
compact subinterval of $(0,T)$ by a nonnegative smooth kernel.
The mollified derivative is nonnegative, and the mollifications converge
uniformly on smaller compact intervals. Their monotonicity passes to $H$;
continuity then includes $0$ and $T$. Consequently
\[
 F(T)\ge e^{-k^2T}F(0).
\]
For $x\in[a,b]$ we have $x(1-x)\ge2/9$, so $y_0\ge2$ there. Hence
\[
 F(0)=\int_a^b(y_0-1)\varphi\,dx\ge\frac2k,
 \qquad F(T)\ge\frac2k e^{-k^2T}>0.
\]
By Lemma~\ref{lem:plateau}, $F(T)=\int_I(y(T)-\bar y)\varphi\,dx$.
Finally, $\|\varphi\|_{L^2(I)}^2=(b-a)/2=1/6$, so Cauchy--Schwarz gives
\eqref{eq:error}.
\end{proof}

The lower bound holds for the same obstacle and initial state at every
finite $T$, independently of the endpoint controls. It tends to zero as
$T\to\infty$ and makes no assertion against approximate or asymptotic
steering. The counterexample uses the freedom to choose an arbitrary smooth
obstacle in Problem 398; it does not decide control questions restricted to
other contact geometries.

\begin{thebibliography}{9}
\bibitem{geshkovski}
B.~Geshkovski, \emph{Control in moving interfaces and deep learning},
PhD dissertation, Universidad Aut\'onoma de Madrid, 2021,
Section 1.5.1, pp.~37--39.
\href{https://cmc.deusto.eus/wp-content/uploads/PhDThesis-bGeshkovski-may2021.pdf}{Institutional full text}.

\bibitem{pighin}
D.~Pighin and E.~Zuazua, Controllability under positivity constraints of
semilinear heat equations, \emph{Mathematical Control and Related Fields}
\textbf{8} (2018), 935--964.
\href{https://doi.org/10.3934/mcrf.2018041}{doi:10.3934/mcrf.2018041}.
\href{https://arxiv.org/abs/1711.07678v4}{arXiv:1711.07678v4}, Section 7.3.

\bibitem{aim398}
AIM problem collection, Problem 398,
\emph{Boundary control to equilibrium for a parabolic obstacle problem},
revision \texttt{c9929805cece044705abaaebe1cebddac1ba11a6}.
\href{https://github.com/MColbrook/AIM/blob/c9929805cece044705abaaebe1cebddac1ba11a6/problems/398-parabolic-obstacle-boundary-control.md}{Pinned statement}.

\bibitem{gerard}
P.~G\'erard, \emph{Distributions and Partial Differential Equations},
lecture notes, December 12, 2019,
Propositions 1.8.2, 1.8.5 and 1.8.7, pp.~35--38.
\href{https://www.imo.universite-paris-saclay.fr/~patrick.gerard/Distributions2019_Chap1,2,3.pdf}{Author's institutional full text}.

\bibitem{hunter}
J.~K.~Hunter, \emph{Notes on Partial Differential Equations},
revised June 18, 2014, Proposition 3.22, pp.~57--58.
\href{https://www.math.ucdavis.edu/~hunter/pdes/ch3.pdf}{Chapter 3}.
\end{thebibliography}

\paragraph{AI assistance.}
OpenAI Codex was used to develop the counterexample and its proof, check
the cited passages, and draft this note. Independent mathematical review
is requested.
\end{document}
